\documentclass[../../../clio.tex]{subfiles}
\begin{document}

\chapter{Repair as Constraint Reconfiguration}
\label{ch:chapter-10-repair-reconfiguration}

\section{Chapter Preface}
Scientific progress repeatedly shows that some failures cannot be solved by searching harder under current assumptions. Sometimes the representational rules themselves must change.

CLIO distinguishes that transition sharply: optimization navigates a fixed admissible space, whereas repair changes the constraint system that defines admissibility.

This chapter formalizes when repair is warranted and how minimal structural change can recover solvability without collapsing semantic coherence.

\section{Derivations and Formal Results}
Let constraint system be \(\constraintSet\) with \(A_y(\constraintSet)=\varnothing\). A repair operation maps constraints by
\[
R:\constraintSet\to\constraintSet'.
\]
Define repair objective
\[
\constraintSet^\star=\arg\min_{\constraintSet'}\left[K(\constraintSet,\constraintSet')+\lambda\,\mathcal I(\constraintSet')\right]
\]
subject to \(A_y(\constraintSet')\neq\varnothing\).

\begin{theorem}[Repair warrant]
If \(A_y(\constraintSet)=\varnothing\) but there exists \(\constraintSet'\) with \(A_y(\constraintSet')\neq\varnothing\), then no state-space optimization under \(\constraintSet\) can solve the reconstruction problem; constraint-level intervention is necessary.
\end{theorem}

\end{document}
